\section{Asymptotics of the balanced observable}\label{c01-c1-asymptotics-of-the-balanced-observable}

The boundary exponent alone does not determine marked bulk mass. We estimate the balanced polygon observable by comparing its color-sum numerator with the squared homogeneous empty component. Their leading exponential factors cancel; control of the remaining powers must be uniform in the complex twist.

\subsection{Color formula and small-cycle gas}\label{c01-c1-color-formula-and-small-cycle-gas}

For the balanced homogeneous case (circumference \(2n\)) the target is the following estimate, with \(C\) uniform on each compact subset of \(|\Im s|<\pi\):
\[
 |X_{n,n}(\kappa=e^s)|\le C\ n^{\,1/6+2\Re(s^2)/(3\pi^2)},\qquad |\Im s|<\pi ,
\]
The normalization is \(X_{n,n}=1\) at \(s=i\pi/2\). Together with positivity the estimate will give the matching exponent at \(s=0\) and analytic control near zero polygon fugacity. The calculation uses (i) a finite color interpolation formula converted to residues; (ii) a one-sided homogeneous normalization estimate for the Pfaffian; (iii) analytic radial periodization of the residues; upper-circle resummation creates hole variables, after which only a small defect gas around an explicit pressure remains. We bound that gas after a Gaussian count summation and recover the nonperiodized value by Cauchy at zero radial nome. The radial period \(L\) is an auxiliary \textbf{spectral} period, not the circumference of the physical SAW cylinder. Matrix and scalar abbreviations will be local to the computations.

Here is a rational product sum for
\[
 \mathcal P=S_N^2 X_{n,m}/\prod_{i<j}x_i x_j
\]
(with the denominator meaning \(\prod_{i<j}(x_i x_j)\)). Assign labels \(a_i\in\{1,0,-1\}\), total sum zero. The type of site \(P={\rm A},{\rm B}\) has sign \(s_P=+,-\) respectively. Take single-site activities \(\lambda_0^P=1\), \(\lambda_+^P\lambda_-^P=\sqrt2-2\) and \(\lambda_+^{\rm A}/\lambda_+^{\rm B}=T\kappa\). Pair factors in order \(P\to Q\) with argument \(r=x_{\rm target}/x_{\rm source}\) and labels \(a,b\) are
\[
 l^{PQ}_{ab}(r)=u^{PQ}_{ab}\,r^{-1-ab}\prod_{j\bmod8}(r-t^j)^{d^{PQ}_{ab}(j)} .
\]
The leading multipliers are \(t^{-s_P b}\) for \(a=0\); for \(a\ne0\) they are \(t^{s_Q a(1-2 b^2)}\). Writing \(e_j\) for unit arrays modulo 8 and * for index negation, the tables in order \(+,0,-\) are
\[
 d^{AA}=\begin{pmatrix}H&J&L\\J^*&H-e_1-e_{-1}&J\\ L^*&J^*&H\end{pmatrix},
\]
\[
 H=e_2+e_3+e_5+e_6,\quad J=e_3+e_5+e_6-e_0,\quad L=e_5+e_6-e_0-e_1,
\]
\[
 d^{AB}=2\begin{pmatrix} e_5+e_6&e_5&0\\e_5&e_4&e_3\\0&e_3&e_2+e_3\end{pmatrix}.
\]
Obtain BB and BA by reversing \textbf{both} labels in AA and AB respectively. Then \(\mathcal P\) is the sum over these neutral assignments of the products of all site and pair factors.

For checking, normalize by the source sign, putting \(g=s_P a,\ h'=s_P b,\ \epsilon=s_P s_Q\). Thus the divisor array is \(D^\epsilon_{g h'}=d^{AA}_{g h'}\) or \(d^{AB}_{g h'}\) for \(\epsilon=+,-\), of total sum \(2+2gh'\). Write \(t^{U^\epsilon(g,h')}\) for the leading multiplier (\(U=-h'\) at \(g=0\), else \(\epsilon g(1-2h'^2)\)). Put \(\tau_k e_j=e_{j+k}\). The tables have \(D^\epsilon_{\epsilon h',\epsilon g}=(D^\epsilon_{g h'})^*\) and \(U^\epsilon(g,h')-U^\epsilon(\epsilon h',\epsilon g)=-\sum_j j D^\epsilon_{g h'}(j)\) modulo 8.

We check the characterization in Proposition~\ref{prop:pv:pairing-interpolation}. Factors are reciprocal on exchanging ends, so order does not matter. Simple poles at equal same-block sites cancel by exchanging unequal labels. At ratio \(t\) the two polar assignments \((s,-s),(0,0)\) in type \(P\) (\(s=s_P\)) cancel including activities. At ratio \(T\) only \((s,-s)\) remains; its local multiplier including activities is \(2-\sqrt2\). At ratio \(t^2\) the possible pairs \((s,0),(0,-s),(s,-s)\) coalesce to their sum with local multiplier twice the single activity. These are the two \(S^2\) reductions with the monomial denominator. To see the spectator factors explicitly put \(B_{ab,c}(r;k)=l^{PQ}_{ac}(r)l^{PQ}_{bc}(r/t^k)\). Substitution in the two tables gives
\[
 \begin{split}
 B_{s,-s,c}(r;1)&=t^{3s c}B_{00,c}(r;1),\\
 B_{s,-s,c}(r;3)&=T^{s c}\frac{(r-t^5)^2(r-t^6)^2}{T r^2},\\
 B_{ab,c}(r;2)&=\frac{(r-t^5)^2}{tr}\ l^{PQ}_{a+b,c}(r/t)
 \begin{cases}1&a+b=\pm1,\\t^{3s c}&(a,b)=(s,-s).\end{cases}
 \end{split}
\]
on the allowed pairs. Indeed for a fixed target label \(c\), abbreviating rows \(D_g=D^\epsilon_{g,sc}\), one reads
\(D_++\tau_1D_-=(1+\tau_1)D_0,\ D_++\tau_3D_-=2(e_5+e_6)\), and \(D_g+\tau_2D_{h'}=2e_5+\tau_1D_{g+h'}\) for \( (g,h')=(+,0),(0,-),(+,-)\); compare the powers and \(U\)'s at infinity. The phases disappear by neutrality. Locally \(l^{PP}_{00}/l^{PP}_{s,-s}\) tends to \(2-\sqrt2\) at \(t\), \(l^{PP}_{s,-s}(T)=-1\), and the three factors at \(t^2\) for the indicated surviving pairs are \(2,2,-(2+\sqrt2)\). At extreme parameters a site of label \(a\) grows with degree \(N-1-a^2\) in either divergent direction by neutrality; only its zero label contributes at top order and the leading phases cancel. This gives polynomiality after restoring \(\prod_{i<j}(x_i x_j)\), the degrees and both extreme decimations.

Finally at the balanced specialization cross factors at \(T\) require the sum of matched labels nonnegative, hence zero at every match. The inter-match factors do not depend on labels: with A label \(a\) matched by \(-a\), and similarly \(b,-b\), they are
\[
 l^{AA}_{ab}(r)^2 l^{AB}_{a,-b}(Tr) l^{AB}_{b,-a}(T/r)
 =r^{-4}\prod_{d\in\{2,3,5,6\}}(r-t^d)^2 ,
\]
as required by the cross factors in the cup reduction of \(S^2\) with the monomial denominator (here \(2D^+_{ab}+(\tau_3+\tau_{-3})D^-_{a,-b}=2H\)). Local multipliers including activities are \((2-\sqrt2)\kappa^a\). This verifies the last interpolation condition and proves the sum formula.

Now specialize to two homogeneous blocks of size \(n\). Factor out the default term (all labels \(s_P\)), which equals
\[
 B_n=(2D_0)^{n^2}D_0^{-2n}\kappa^n,\qquad D_0=2+\sqrt2 .
\]
The remaining gas has objects in each type of length one (denoted E) or two (D); the total virtual length \(M\) is the same in both types. An E object changes a site label from its default $s_P$ to zero; a D object changes it to $-s_P$. Object length counts virtual integration variables, rather than visited vertices of a walk. Integrate each top variable \(z\) around a small positive circle surrounding 1 with measure \(dz/(2\pi i z)\); the virtual variables of an object are \(z\) alone or \((z,z/t)\) respectively. Each virtual variable carries \(h^n\) evaluated there, and distinct objects carry all pairwise interactions between their virtual variables, as follows:
\[
 h(r)=\frac{(r+t)^2}{(r-1)(r-t^2)},\quad
 K_s(r)=\frac{(r-1)^2(1+r^2)}{t r(r-t)(r-1/t)},\quad
 K_x(r)=\frac{t r(1+r)^2}{(1+\sqrt2 r+r^2)^2}
\]
(pairwise argument = ratio, same and different colors respectively). Sum with inverse factorials of the numbers of indistinguishable objects and activities
\[
 A_P=\frac{1}{a_0\lambda^P_{s_P}},\quad a_0=\frac{t^2}{1-t^2},\qquad
 B_P=-A_P^2\,\mathop{\rm Res}_{r=t}\big(K_s(r)/r\big)
\]
for E,D respectively. This gives \(\mathcal P/B_n\).

For the residue check first separate base sites slightly, preserving equal block sizes and dividing by their corresponding default term, using instead of \(h^n\) products at virtual/site ratios of \(p_s(r)=t^2 r/[(r-1)(r-t^2)]\) on same-color sites and \(p_x(r)=t^6(r-t^5)^2/r\) on the others. Objects now select distinct sites of their color (coincident tops have same-color zero of order 2,1,2 for EE,ED,DD in the combined kernels). Length \(i\) means new label \(s_P(1-i)\). By the pair tables, against a still-default site the change ratio equals the string of \(p\)'s times \(t^{3s_Ps_Q 1_{i=2}}\). Between two changed sites the correction (on top of their two single-site changes) is the string of \(K\)'s times
\(t^{-3s_Ps_Q(i 1_{j=2}+j 1_{i=2})}\); these identities follow just by multiplying linear factors. The constants thus give net \(t^3\) per D: for a length-2 site in type \(P\) the assigned exponent is \(3s_P\sum s_Q(1-j)\) over other sites of lengths \(j\) (taking \(j=0\) if unchanged), which equals 3 by neutrality. The extra self-site residue factors are \(a_0\), or \(a_0 p_s(1/t)\) for a D, consistent with \(B_P=t^3\lambda^P_{-s_P}/(\lambda^P_{s_P} a_0 p_s(1/t))\). No other poles intervene on these small contours. This proves the integral formula (one can truncate to \(M\le 2n\)).

For instance use the ratio (top at type \(P\))/(top at \(Q\)), with \(\epsilon=s_Ps_Q,\ g=1-j,\ h'=\epsilon(1-i)\) normalized now at source \(Q\). The correction divisor read from the tables is
\[
 D^\epsilon_{g h'}+D^\epsilon_{1,\epsilon}-D^\epsilon_{1,h'}-D^\epsilon_{g,\epsilon}
 =\sum_{u=0}^{i-1}\sum_{v=0}^{j-1}\tau_{u-v} K_\epsilon^\#,\qquad
 (K_+^\#,K_-^\#)=(2e_0+e_2+e_6-e_1-e_7,\ 2e_4-2e_3-2e_5),
\]
and the analogous difference of four \(U\)'s modulo 8 is \(\epsilon[-ij+ij(j-i)/2-3(i1_{j=2}+j1_{i=2})]\), giving the stated phase including the shifts in the kernel string.

\subsection{Homogeneous Pfaffian normalization}\label{c01-c1-homogeneous-pfaffian-normalization}

\begin{proposition}[Pfaffian denominator]\label{prop:pv:denominator}
The homogeneous empty component satisfies (\(D_0=2+\sqrt2,\ \delta=3/8\))
\[
 S_{2n}(1,\ldots,1)\ge c\,D_0^{-n}(4D_0^2/9)^{n^2} n^{5/48}.
\]
\end{proposition}

\begin{proof}
The lower bound includes bounded error in its logarithm, because a larger normalization error would affect the final power. The kernel at variables \(e^{2ia}, e^{2ib}\) equals \(i\sin(2(a-b))/(\cos(2(a-b))+1/\sqrt2)\), represented by \(i\,\mathrm{pv}\int e^{(a-b)p}W(p)\,dp/p\) as in the elementary transforms above, with
\[
 W=\frac{p\cosh(\lambda p)}{2\sinh(4\lambda p)}
 =W_0 e^{f(p)},\quad W_0=\frac{p}{4\sinh(\delta\pi p)},\quad f(p)=\log(1+1/(2\cosh(\pi p/4))).
\]
Taking the homogeneous Taylor/Vandermonde limit in the Pfaffian gives
\[
 S_{2n}(1)= (D_0/2)^{\binom{2n}2}\frac{\det[\int p^{2j+2k} W(p)\,dp]_{j,k=0}^{n-1}}{\prod_{j=0}^{2n-1}j!}.
\]
By multilinearity the Pfaffian over Taylor coefficients uses indices \(0,1,\ldots,2n-1\) (leading alternating order); paired evens and odds give this determinant, signs cancelling against the prefactor (\(i^{\binom{2n}2}(-i)^n=1\)). For \(W_0\) instead, using \(y=\delta p, C=\int W_0=1/(8\delta^2)\), even coefficients of
\[
 (1-v)^{-1-iy}(1+v)^{-1+iy}
\]
are real orthogonal polynomials with squared norms \(C(m+1)\) at index \(m=0,2,\ldots\). Just integrate two generating functions: the normalized push \(2y\,dy/\sinh(\pi y)\) has Fourier transform \(1/\cosh^2(a/2)\), giving \((1+vu)^{-2}\). Thus the monic squared norms are \(C m!(m+1)!/(2\delta)^{2m}\), and the formula with \(W_0\) gives precisely \(D_0^{-n}(4D_0^2/9)^{n^2}\).

Log of the moment determinant ratio is bounded below by the trace of multiplication by \(f\) in the reference even-polynomial space, by scalar Jensen in an orthonormal eigenbasis of the positive Gram matrix. We give details to order \(O(1)\). For real \(|y|\le c_1\log m\) at large even \(m\) the unnormalized generating coefficient equals
\[
 2\Re\left[2^{-1+iy} m^{iy} H(y)(1+O((1+y^2)/m))\right]+O(m^{-1/2}),\quad
 H(y)=\lim_{k\to\infty} k^{-iy}\prod_{j=1}^k(1+iy/j).
\]
Choose \(c_1>0\) sufficiently small. Indeed subtract from the generating function the two singular terms \(2^{-1+iy}(1-v)^{-1-iy},2^{-1-iy}(1+v)^{-1+iy}\). The derivative of the remainder in the disk is bounded by \(C(1+y^2)e^{C|y|}/|1-v^2|\) by a first-order estimate of the regular factor at each end. On radius \(1-1/m\) the angular integral of \(1/|1-v^2|\) is \(O(\log m)\), proving the additive error via coefficient extraction of the derivative. The multiplicative error follows by the convergent logarithmic products (tails after subtracting first-order terms are \(O(y^2/m)\)). Also \(|H(y)|^2=\sinh(\pi y)/(\pi y)\).

Drop the range \(|y|>c_1\log m\) by positivity in the trace lower bound. The remaining integral per even index is at least
\[
 \frac{1}{\pi(m+1)}\int_{\mathbb R} f(y/\delta)
 \left[1+\Re\big((2m)^{2iy}H(y)/\overline{H(y)}\big)\right]dy
\]
up to summable errors: integrate the square of the coefficient against \(2y f(y/\delta)\,dy/[(m+1)\sinh(\pi y)]\); \(|H|^2\) cancels its weight, and \(f\) decays exponentially (also to restore the omitted range on the right). The oscillatory integrals are \(O((\log m)^{-2})\) before the prefactor, by parts twice; the first two logarithmic derivatives of the phase factor \(H/\bar H\) have polynomial growth by its product. Finally \(\int f(p)dp=5\pi/9\), by expanding \(\log(1-e^{-3u})-\log(1-e^{-u})-\log(1+e^{-2u})\), \(u=\pi|p|/4\) (using \(\sum k^{-2}=\pi^2/6\), e.g. by Parseval on the linear function). The even-index sum now gives the claimed lower bound.

\end{proof}

The denominator is now fixed. We return to the numerator integral and deform it without changing its physical value at zero radial nome.

\subsection{Radial periodization of the small-cycle integral}\label{c01-c1-radial-periodization-of-the-small-cycle-integral}

Write \(\alpha=\pi/4,\ \kappa=e^s\), and use coordinates \(x=\log(z/t)\), periods \(L,2\pi i\), radial nome \(\mathfrak q=e^{-L/2}\), and \(P=\mathfrak q^2\). Define
\[
 S(x)=2\sinh(x/2)\prod_{a\ge1}(1-P^a e^x)(1-P^a e^{-x}).
\]
Use kernels \(\prod_j S(x-i\alpha j)^{d_j}\) specified by these divisors (actual integer shifts, so \(4,-4\) denote two factors):
\[
 h:\ d_{\pm1}=-1,\ d_{\pm4}=1;\quad
 K_s:\ d_0=2,\ d_{\pm1}=-1,\ d_{\pm2}=1;\quad
 K_x:\ d_{\pm3}=-2,\ d_{\pm4}=1.
\]
They are for the external field in \(x\) and for pair differences respectively. At \(P=0\) these agree with the small-cycle kernels with multipliers \(t^{-1},t\) stripped from the pair kernels; these phases would contribute \(t^{M+\#D}\) in total. Set \(C_L^2=\lim_{x\to0}K_s(x)/x^2\) (square root positive at \(P=0\)), \(k_L={\rm Res}_{x=i\alpha}K_s(x)\). Use E activity \(a_L=-i C_L e^{-s/2}\) in each type and D activity \(-k_L a_L^2\) (incorporating the phases at \(P=0\); indeed \(C_\infty^2=D_0\) and \(t A_A A_B=-D_0/\kappa\) there). Multiply the integrand by
\[
 \Theta(C)^2,\quad C=\exp(\sum_\nu q_\nu x_\nu),\quad
 \Theta(C)=\sum_{b\in\mathbb Z} e^{-Lb^2/2} C^b
\]
where \(\nu\) ranges over virtual variables including both constituents of D, and \(q_\nu\) is their color sign. D constituents are now \(x,x-i\alpha\). Integrands are elliptic in each object's top coordinate, by neutrality and the shift formula \(S(x+L)=-e^{x+L/2}S(x)\) (the multiplier of the pair kernels moving one object of charge \(p_i\) is \(C^{-2p_i}\exp(-p_i^2L)\), canceled by \(\Theta(C)^2\)).

Denote by \(\mathcal N_n(\mathfrak q,s)\) the resulting small-cycle gas around tops \(x=-i\alpha\), summed over \(0\le M\le2n\) with its factorial measures and deformed activities, and multiplied by \(B_n\). This includes the factor \(\kappa^n\) in \(B_n\). The numerator \(\mathcal N_n\) is analytic for small \(\mathfrak q\); at zero nome the homogeneous monomial denominator in \(\mathcal P\) is one, so
\[
 \begin{aligned}
 \mathcal N_n(0,s)&=\mathcal P(1,\ldots,1;\kappa=e^s)\\
 &=S_{2n}(1,\ldots,1)^2 X_{n,n}(\kappa=e^s).
 \end{aligned}
\]
Even if the small-cycle sum is extended, any fixed \(M>2n\) contributes zero: resolve \(h(x)^n\) into a product with slightly displaced fields and use distinct top poles as before, then pass to the limit.

Here a single sufficiently small contour size and disk of analyticity in \(\mathfrak q\) work independently of \(n\). For top differences near zero the same-color string products have orders \(2,1,2\) for EE, ED, DD, and the opposite-color ones have no poles; other real-period images stay away. So after splitting the external poles each color selects distinct top poles (at most \(n\), length at most 2 each). Products defining \(S,\Theta\) converge normally where needed, and \(C_L,k_L\) are analytic deformations of their values at zero radial nome.

We describe term transformations first at fixed \(M\), large real \(L\). Write a radial cycle of shift \(\theta\) as \(x=r-i\alpha\theta,\ -L/2<r<L/2\), positive towards the right. Expand each small cycle as lower minus upper (shifts \(1\pm\epsilon_0\), tiny \(\epsilon_0\)). Move upper E's up to shift 0. The pole crossed against each still-lower E of the same color is simple, producing a two-string with coefficient \(-k_L a_L^2\) (minus sign of the moving cycle). These cancel all upper D contours, which have opposite coefficient. More explicitly one can move the E contours successively; merged pairs have top near shift 0 and can move to shift \(1-\epsilon_0\). They have no further poles against remaining E or D's in this step (same-color nearby poles of string products cancel against zeros), hence the binomial expansion over disjoint mergings gives the stated cancellation with the factorials. Cycles can be infinitesimally separated during moves. Move both lower contours together to shifts \(2+\ell\) for remaining E (now called \(\mathrm L\)) and \(2+d\) at D tops, with \(\ell=1/2, d=-1/8\). No pole is crossed (opposite-color pair distances from the upper contour remain below 3; lower top relative shifts stay smaller than 1). The other E's are now at shift 0 (\(\mathrm U\)) with reversed sign.

In this deformation the small-cycle expansion holds by periodicity since all top contours in the initial small band have no pair poles. For each choice of \(r'\) disjoint mergings in one color, the two E factorials after summing matchings leave factorials for the two residual E counts and an extra \(1/r'!\). The residue equation is \(x_{\rm moving}=x_{\rm lower}+i\alpha\), so after moving the created contours next to the upper D contours, a fixed total \(d'\) there has relative coefficient \(\sum_{r'=0}^{d'}(-1)^{d'-r'}/[r'!(d'-r')!]\) times the D activities. Hence only \(d'=0\) remains.

For instance an E approaching a D with their virtual shift differences at \((1,2)\), \((-1,0)\) or \((0,1)\) sees no same-color pole in the product; two D's likewise have no pole for top differences \(0,\pm1\). Newly created tops can go from \(\epsilon_0\) to \(1-\epsilon_0\) with bottoms strictly beyond 1. At the last move take linear interpolation of both lower top shifts: no U--lower constituent distance reaches 3 or recrosses 1, and the L,D top relative shift has size \(<1\).

Here are explicit zero and nonzero mode conventions. Put \(\varepsilon=2\pi\alpha/L,\ \phi=2\pi r/L,\ y=e^{i\phi}\). For a pair of constituents with kernel divisor \(d_j\), shift difference \(\theta\) (target minus source), \(p=q_{\rm source}q_{\rm target}=\sum d_j/2\), use
\[
 W(v)=\sum_j d_j\frac{\exp[v((j+\theta)\bmod8-4)]}{2\sinh4v}
 =\frac{p}{4v}+W^{[0]}+W^{[1]}v+O(v^2),\quad
 k=\sum_j d_j(1/2+\lfloor(j+\theta)/8\rfloor)
\]
taking half averages of the two branches at a wall. For whole objects sum the arrays over constituents, obtaining \(W_{ij}\) when source and target are lower objects (also use copies at identical data for \(i=j\)), a two-color column \(J_i^+\) from \(\mathrm U\) to lower \(i\) (reversed = \(J_i^-\)), a symmetric matrix \(A\) between U variables, and \(H\mathbf1\) (for U) or \(h_i^+\) (for lower, reversed = \(h_i^-\)) from the external field once (\(p=0\)). In particular \(H=(1-\cosh3v)/\sinh4v\), \(A_{ss},A_{sx}\) (diagonal and off-diagonal) equal \((\cosh4v-\cosh3v+\cosh2v)/\sinh4v,(1-2\cosh v)/\sinh4v\), both eigenvalues of \(A\) positive.

These conventions follow from expanding the Gaussian-modified kernel (multiply by \(\exp[-p(x_j-x_i)^2/L]\)). It has winding phase \((-y_j/y_i)^k\), mean log \(p L/6+p\alpha^2\theta^2/L-\varepsilon W^{[1]}\) apart from phase, and log coefficient \(-W(m\varepsilon)/m\) at \((y_j/y_i)^m\), \(m>0\), with reversed shifts at negative frequency. Winding powers here can be grouped to integer powers for whole objects (off the walls this follows from even total divisor order per kernel; at the final shifts the same-contour EE walls have even order, and for DD the \(+1,-1\) differences of constituents enter together along with the two zero differences). One direct derivation uses, for each shifted sine at \(r-i\alpha u\),
\[
 S=\mathrm{sgn}(r)e^{(|r|-i\alpha u\,\mathrm{sgn}(r))/2}
 \prod_{l\in\mathbb Z}(1-e^{-|r-lL|} e^{i\alpha u\,\mathrm{sgn}(r-lL)})
 \qquad (|r|\le L/2).
\]
Signs cancel; the log jump at 0 after summing and Gaussian compensation is \(-2\pi i k\), giving the indicated phase. The mean of the product log per shifted factor is \(-2\pi^2 B_2(\{u/8\})/L\), \(B_2(w)=w^2-w+1/6\). At nonzero frequency expand the logs; integrating on the two half-lines gives sums of \(1/(a^2+\omega^2)\) against \(\cos(a\alpha u)\) and \((\omega/a)\sin(a\alpha u)\), \(\omega=2\pi m/L\). Combining with the exponential and winding terms yields the formula for \(W\) (equivalently use the Fourier series of \(e^{\omega(x-\pi)}/(2\sinh\pi\omega)\) on \(0<x<2\pi\)). At zeros take log-moduli singularities by radial limits.

In detail, for \(\beta=\alpha u\bmod 2\pi\), per unit divisor the positive log coefficient after winding removal is \(-1/L\) times
\[
 \frac1{\omega^2}+\frac{\beta-\pi}{\omega}
 +2\sum_{a\ge1}\frac{\cos(a\beta)+(\omega/a)\sin(a\beta)}{a^2+\omega^2}
 =\frac{2\pi}{\omega}\frac{e^{\omega(\beta-\pi)}}{2\sinh(\pi\omega)}
\]
(average endpoints), using \(i(\phi-\pi\,\mathrm{sgn}(r))\) as the two-half log of \(-y\) and symmetry of the shifts to distribute the Gaussian terms.

Poisson-transform the two theta factors (dual indices \(b_1,b_2\), sum \(b\)). Let \(p_i=\sum_{\nu\in i}q_\nu\), \(z_0=\sum_\nu q_\nu\theta_\nu\). Let \(u\) be the column of U counts (length \(M\) minus lower virtual counts in each type); let \(e\) be the counts of lower objects (D counted once). Apart from activities, measures, powers \(e^{-sM}\), and lower-object self corrections specified below, the scalar mean is \(2\pi/L\) times
\[
 \exp\left\{-\frac{\varepsilon}2\left[
 (z_0+4b)^2/4+4(b_1-b_2)^2
 +[v]\{u^t A u+2u^t(nH\mathbf1+\sum J_i^+)
 +\sum_{ij} W_{ij}+2n\sum h_i^+\}\right]\right\},
\]
where \([v]\) takes the coefficient of \(v^1\) at zero. Indeed the raw-pair Gaussian compensation sums to \((-(\sum q_\nu x_\nu)^2+\#D\,\alpha^2)/L\) by neutrality, the square canceling the modular exponential. Pair mean offsets sum to \(\alpha^2(-z_0^2-\#D)/L\) (the internal D pairs are omitted). The \([v]\) terms use the full ordered square of all objects and require diagonal corrections below. Phases outside the upper expectation are
\[
 \prod_i y_i^{-b p_i}(-y_i)^{u_{c(i)}-n |p_i|}
 \prod_{i<j}(-y_j/y_i)^{k_{ij}}
\]
where \(c(i)\) denotes color. Upper means and smooth self factors convert upper cycles exactly to independent circular unitary measures in each color (Vandermonde modulus squared with normalized Haar circles and factorials). The inserted winding power on each upper \(y\) is \(-e_{c}-bq_c\). Additionally insert, in the exponent for each frequency \(m>0,\ v=m\varepsilon\), with \(T_\pm\) the columns of sums of upper \(y^{\pm m}\), the quadratic term
\[
 m^{-1}\big[T_-^t(1-A)T_+
 -T_-^t(nH\mathbf1+\sum_i J_i^+y_i^m)
 -(nH\mathbf1+\sum_i J_i^- y_i^{-m})^t T_+\big].
\]
The remaining nonzero-mode log is
\[
 -m^{-1}\big[\sum_{ij}W_{ij} y_j^m y_i^{-m}
 -\sum_i W_{ii}(+\infty)(1-e^{-v})
 +n\sum_i(h_i^+y_i^m+h_i^- y_i^{-m})\big]
\]
per mode.

Indeed from U to a lower object the winding \(k\) is 1 in the same color and 0 in the other, and for its external field is \(-|p_i|\). Fourier sums with coincident cycles are understood away from coincidences or with the pure log singularities factored. Including self terms requires the factor \(\exp(\Gamma_i)\) per lower object (note \(W_{ii}^{[0]}=0\) by reversal symmetry), where
\[
 \Gamma_i=-p_i^2 L/12+\varepsilon W_{ii}^{[1]}/2+
 \sum_{m>0}[W_{ii}(m\varepsilon)-W_{ii}(+\infty)(1-e^{-m\varepsilon})]/m=O(1)
\]
uniformly at large \(\Re L\), bounded \(\Im L\) (subtract \(p_i^2/(4m\varepsilon)\) in the series to cancel the linear growth). At U instead subtract 1 from the same-color coefficient at every frequency for the Haar separation: the resulting self log correction is exactly \(-\log(LC_L/(2\pi))\) by the limiting same-color kernel divided by \(|1-e^{i\phi}|^2\). Thus it cancels the signed measure conversion with \(a_L\), modulo the length fugacity already counted. Lower objects still have factorial measures and activities with \(e^{-sM}\) removed. These fixed-\(M\) identities continue to complex \(L\) (radial cycles along the period) near large positive real parts with bounded imaginary parts.

In this bookkeeping \(\sum_{i<j}p_i p_j L/6=-\sum_i p_i^2 L/12\) over distinct whole objects (including U). Explicitly for real period the U same-color compensated kernel at shift zero divided by \(|1-e^{i\phi}|^2\) tends at coincidence to \((LC_L/(2\pi))^2\). Its logged limit by the mean and Fourier coefficients above is
\[
 L/6-\varepsilon A_{ss}^{[1]}-2\sum_{m>0}(A_{ss}(m\varepsilon)-1)/m .
\]
Negative half of this is the self correction for writing the smooth residual U interactions with both indices independently summed. The kernel at this alignment is positive on the real cycle away from its zeros, so there is no extra phase in the conversion.

\subsection{Upper-edge resummation}\label{c01-c1-upper-edge-resummation}

The upper contours still contain an arbitrarily large number of points. The following circle determinant expansion sums them first and replaces them by holes on two separated contours. The dependence of its sign on the winding and circle size must be retained: it is what later permits the Gaussian sum over the original count \(M\).

The zero-winding case is closely related to the Toeplitz--Fredholm identity of Borodin and Okounkov~\cite[Theorem~1]{BorodinOkounkov2000}. The calculation below retains the winding and its size-dependent sign; its analytic interchange with the Hubbard fields is justified separately.

Here is the winding version of the expansion. Take potentials \(V_+(y),V_-(y)\) with strictly positive and negative powers respectively, analytic on a neighborhood of a closed annulus spanning the unit circle. For a circular unitary expectation of size \(u\) insert winding \(y^{-d}\) and \(\exp(V_++V_-)\) on each variable. Expand in additional points called holes, with signs \(\tau=+1\) (outer circle \(>1\), O) or \(-1\) (inner \(<1\), I), counts \(r,l\) with \(r-l=d\). Use normalized positive circle measures divided by factorials. The scalar prefactor is
\[
 (-1)^{ud}\exp\sum_{m>0}m [y^m]V_+ [y^{-m}]V_- .
\]
Hole factors are
\[
 \prod Y_h^{-u\tau_h}\exp\sum_h\tau_h(V_-(Y_h)-V_+(Y_h))\ 
 \prod_{h<h'} E(Y_h,Y_{h'})^{\tau_h\tau_{h'}},\qquad E(a,b)=(a-b)^2/(ab)
\]
times a sign depending on \(r,l,d\) only (not \(u\)); only this independence is needed. Circles can be taken wherever the series potentials converge in the indicated annulus.

We give proof details to specify both convergence and the dependence of phases. It suffices first to use polynomial positive and negative times, as formal series. For \(d\ge0\), expand the two exponentials against Vandermonde determinants by minors and integrate: one gets
\(\sum_\lambda s_{\lambda+(d^u)}(V_+)s_\lambda(V_-)\) with \(\ell(\lambda)\le u\). Here \(s_\lambda(t)=\det[h_{\lambda_i-i+j}]\), \(h(z)=\exp(\sum_{k>0}t_k z^k)\) for the respective coefficient list. Transposing partitions gives \((-1)^{ud}\sum_{\nu_1\le u}s_{(u^d,\nu)}(-V_+)s_\nu(-V_-)\). Indeed the partition determinant is a Toeplitz minor on columns \(u-i\), rows \(u-i+\lambda_i\); the complementary inverse minor on a sufficiently large finite index interval starting at 0 uses the complement \(u+j-1-\lambda^t_j\) and \(h^{-1}\), giving the transpose and sign.

Explicitly \(\det[y_j^{u-i}]\prod_j h(y_j)\) expands by matrix multiplication into determinants with rows \(y_j^{u-i+\mu_i}\) for partitions \(\mu\) (padded to size \(u\)), with precisely these minor coefficients. The reciprocal-variable determinant expands likewise. Their product integrates with \(\prod_j y_j^{-d}\) and the circle factorial measures by orthogonality of monomials, requiring positive-list indices to be the negative-list indices shifted up by \(d\); this gives the asserted first sum. In the transposed sum \((u^d,\nu)\) means prepend \(d\) parts of size \(u\).

The charged-wedge calculation is the classical boson--fermion correspondence~\cite{Frenkel1981}; see the Clifford-operator and Schur-function treatment in \cite[Theorems~8.1 and~8.3]{FrenkelJingWang2000}. We give its coefficient calculation to fix the charge convention and the projection onto permitted occupied levels, before taking any analytic specialization.

Use wedges with decreasing occupied positions \(c+\lambda_i-i\), charge \(c\), vacancy of the partition denoted \(|c\rangle\). Insert/remove level \(j\) by \(\psi_j,\psi^*_j\) (sign by exterior order). The currents \(A_k=\sum_j\psi_{j-k}\psi_j^*\ (k\ne0)\) satisfy
\[
 [A_k,\psi_j]=\psi_{j-k},\quad [A_k,\psi^*_j]=-\psi^*_{j+k},\quad
 [A_k,A_h]=k\delta_{h,-k}.
\]
Indeed terms act finitely on each state, and the last bracket commutes with all insertions/removals so is scalar (commute first with occupation projectors); evaluate on the vacancy, using the \(k\) possible moves up by \(k>0\). The partition amplitudes of \(\Gamma_-(t)|c\rangle\) with \(\Gamma_\pm(t)=\exp(\sum_{k>0}t_k A_{\pm k})\) are the minors above by exterior multiplication with the positive shift generating function \(h\) (can truncate the sea deep below all possible movements at any fixed weighted degree). Thus after removing the transpose sign our sum is
\[
 \langle d|\Gamma_+(-V_+)\psi_{u+d-1}\cdots\psi_u
 \prod_{j\ge u+d}(1-\psi_j\psi_j^*)\Gamma_-(-V_-)|0\rangle .
\]
Expand the projectors (finitely effective in each weighted degree) and use generating currents \(\psi(w)=\sum\psi_k w^k,\ \psi^*(v)=\sum\psi_k^* v^{-k}\), with creations outside, annihilations inside. Commuting the \(\Gamma\)'s past the currents gives precisely the exponential factors stated. The vacancy expectation with \(r=d+l\) creations and \(l\) annihilations (ordered separately, signs from reordering depending only on counts) is up to sign the determinant with columns \(w_i^a\ (0\le a<d), v_j/(w_i-v_j)\), by summing over distinct removed negative levels and creating these plus \(0,\ldots,d-1\). The mode extractions, symmetrized using antisymmetry and summed by minors over the projection levels, give a second determinant with columns \(w_i^{-a},(v_j/w_i)^d/(1-v_j/w_i)\), apart from \(\prod_h Y_h^{-u\tau_h}\) and inverse factorials. Indeed the fixed creations use levels \(u+a\), and the \(l\) projector pairs distinct levels \(u+d+n',\ n'\ge0\). After the \(Y_h^{-u\tau_h}\) extraction, antisymmetrization in the \(w\)'s gives columns \(w_i^{-a},w_i^{-d-n'}\), in the \(v\)'s the minor on powers \(d+n'\). Sum over subsets of the \(n'\)'s (listed in order) by determinant multiplication using \(|v_j|<|w_i|\); the antisymmetrizations divide by \(r!\,l!\), with ordering signs depending only on counts. Subtract polynomial columns; the product of the two Cauchy--Vandermonde determinants gives exactly the pair factors up to sign independent of \(u\). All mode sums at fixed counts converge on the separated radii. The series thus obtained converges absolutely for deterministic analytic potentials (Hadamard bounds \(r^r C^{r}\) for the two determinants at fixed \(d\), absorbed by the factorials), giving the analytic identity. Negative \(d\) follows by inverting circles and interchanging the signs.

To elaborate, in the second determinant use the polynomial columns and the finite geometric sum so its non-polynomial columns become \(1/(1-v_j/w_i)\); thus it is a Cauchy--Vandermonde in the reciprocal variables. The product with the first determinant has Vandermonde squares divided by \(\prod(w_i-v_j)^2\), and the remaining powers per \(w_i,v_j\) are respectively \(l-r+1,r-l+1\), exactly as in the \(E\) factors. Projector subset signs, antisymmetrization signs and Cauchy--Vandermonde signs here do not depend on \(u\). For finitely many times the absolute bound above is uniform on their compact sets; Taylor coefficients agree formally (only finitely many projector subsets can act at fixed weighted degree before commuting). Then approximate analytic potentials by truncations on separated circles.

Now use Hubbard linearization of the upper nonzero modes in the preceding integral. Write \(B^\pm=nH\mathbf1+\sum_i J_i^\pm y_i^{\pm m}\) at positive frequency \(m\). In the color basis jointly use coefficient columns
\[
 V_{\pm m}=-B^\mp/m+\sqrt{1-A}\,\xi_{\pm,m}/\sqrt m ,
\]
where the columns \(\xi_{+,m}=\overline{\xi_{-,m}}\) are independent standard circular complex Gaussians across modes (matrix square root in the fixed real orthogonal eigenbasis). Their average reproduces the quadratic term. Apply the winding expansion in each color, with its upper size and its winding \(d_c=e_c+b q_c\) (distinct from lower contour shift notation). Set \(Y_h=e^{\varepsilon h}y_h\), where the shift \(h\) is initially tiny positive on O, negative on I. Let \(c^\pm\) collect \(\sum \tau Y_h^{\pm m}\) by color. Averaging the exponential factors from the expansion gives exactly
\[
 \det A(v)^{-1}\exp\left( m^{-1}
 \{(B^--c^-)^t A(v)^{-1}(B^++c^+)+ (c^-)^t c^+\}\right)
\]
per mode. Hole pair factors \(E^{\tau\tau'}\) apply within each color. Their winding/positive-modulus separation uses mean log \(\varepsilon|h-h'|\tau\tau'\), phase power \(-\mathrm{sgn}(h-h')\tau\tau'\) on \(y_{h'}/y_h\) (up to fixed signs). Thus the resulting nonzero Fourier hole-hole kernel, including the displayed Gaussian terms, is \(\tau\tau'e^{(h'-h)v}(A^{-1}+\mathrm{sgn}(h-h')1)\), with self subtractions for the original circle pair singularities. Hole-lower kernel in forward direction is \(\tau e^{-hv} A^{-1}J_i^+\), in reverse \(-\tau (J_i^-)^t A^{-1} e^{hv}\). The two factors \((-1)^{u_cd_c}\) combine with the preceding lower winding powers to leave no count sign depending on \(M\) since \(u=M\mathbf1-\mathrm{(lower\ virtual\ counts)}\).

For the Gaussian/series interchange at fixed \(L\) real and fixed discrete data before introducing holes, the random potentials are analytic on a small annulus (exponentially decaying high modes). Moreover the series is integrable absolutely. Choose shifts \(\pm h_0\) sufficiently tiny inside this annulus. Integrating the modulus against the Gaussians uses the positive stiffness \(A(v)\) after the Szeg\H{o} prefactor. Apart from one-body terms (bounded per hole at fixed data), the absolute real nonzero-mode Hessian after including Cauchy factors and Gaussian gains, conjugated by charge signs \(\tau\), is, per color eigenvalue \(a\) of \(A\),
\[
 e^{-|h-h'|v}-(a^{-1}-1)_+\sinh(hv)\sinh(h'v)
 -(1-a^{-1})_+\cosh(hv)\cosh(h'v).
\]
Indeed the linear real Gaussian sources from holes have coefficients proportional to \(2\sinh(hv)\Re(\xi y^m)\) for a real-coupled mode, \(2\cosh(hv)\Im(\xi y^m)\) for an imaginary-coupled one. The matrix on the two shift classes is positive definite when \(\tanh(h_0v)<a<\coth(h_0v)\). This holds for tiny fixed \(h_0>0\) by the explicit \(A\) (eigenvalues asymptotic to \((9/8)v,1/(2v)\) at zero and \(1+O(e^{-v})\) at infinity). At fixed \(\varepsilon\), subtract \((1-e^{-cv})1\) for very small \(c>0\), leaving positive smooth Fourier forms whose missing diagonal terms cost at most constants per hole. The subtracted kernels have pairwise gain at most \(O(c\varepsilon)\) per pair (sum the logarithms), which is absorbed by the favorable Cauchy mean \(-2\varepsilon h_0\sum_c r_c l_c\) up to constants since count differences are fixed. The factorial measures now ensure absolute convergence.

Here the small subtraction is the identity on both color and shift indices. Its distinct-point terms only concern identical shifts in one color (hence equal charge signs), with pair gain bounded by \(2\log\sup_{|z|=1}|(1-z)/(1-e^{-c\varepsilon}z)|\le c\varepsilon\). It preserves positivity if \(c\) is small: at high modes the Hessian tends exponentially to identity, and on the rest \(v\ge\varepsilon\) use strict positivity. Completing the residual (positive) ordered sums with their diagonals then only pays \(\sum_{m>0}e^{-cm\varepsilon}/m\) per hole: the Gaussian gains already have their diagonals, and the original missing Cauchy diagonal coefficient 1 was split into \(1-e^{-cv}\) (kept on distinct points only) and \(e^{-cv}\). Also \(\binom{r_c}{2}+\binom{l_c}{2}\le r_c l_c+O(1)\) at fixed count difference for the two same-shift classes in each color (the subtraction constant can be chosen below \(2h_0\)). Gaussian integrations can also be truncated in modes first (for each fixed number of points the high coefficients are exponentially small, giving integrable domination of that tail).

The absolute bounds work uniformly over the lower radial positions here at fixed preceding counts and indices. With Gaussian cutoffs the high-mode gain terms in the real Hessian simply drop out, preserving the bound, and the Gaussian cross terms with the deterministic sources are bounded by a constant per hole by high-mode decay on the tiny annulus. For each fixed number of holes the Gaussian formula passes to the infinite product as well (summable coefficient tails; one can take absolute moments of order slightly above 1 using strict Gaussian stiffness).

\subsection{The residual kernels}\label{c01-c1-the-residual-kernels}

The circle determinant has eliminated the \(\mathrm U\) variables. We first carry out the remaining transformations for a fixed nonnegative integer \(M\) and a large real period \(L\). Here \(M\) is the common total number of virtual integration variables per color in the original small-cycle gas, including variables assigned to lower objects. It is unrelated to the visited-vertex length of a walk. If \(g_c\) is the number of lower virtual variables of color \(c\), the eliminated upper ensemble had size \(u_c=M-g_c\).

We now give the remaining interactions as a matrix on a fixed list of species. The exact one-object measures and factorials are retained, together with the implicit constant phase signs. The one-object prefactors obey the stated uniform bounds and are independent of \(M\) at fixed defect data; the implicit phases have modulus one and depend only on the discrete configuration. The position-dependent phases remain explicit below. Section~\ref{c01-c1-estimates-and-summation-of-the-defect-gas} will justify absolute summability of the assembled expression, complex-period continuation, and the subsequent sum over \(M\).

The hole sign \(\tau=1,-1\) refers to the outer or inner contour of an eliminated upper ensemble. Both signs occur in each gas color, whose sign is \(q\). In the moves below we continue the Fourier coefficients and their winding powers as functions of the offsets, re-expanding a pure singular factor when its Fourier branch changes. The upper circles have already been integrated out.

After integrating the Hubbard fields move hole contours O to \(h=3/8\); then move I to \(h=-1/4\). In the second move the pole with a D of the same color creates an object \(\mathrm P=(\mathrm D,\mathrm I)\) at equality of their radial variables, with \(d=h\) (here \(2+d\) is the lower top shift of D). Put free D at \(d=1/8\) and P at \(d=h=0\). Their unshifted radial variables within the period are denoted \(x_i\) in what follows (so \(y_i=\exp(2\pi i x_i/L)=e^{i\phi_i}\), and e.g. a hole has \(Y=e^{\varepsilon h}y_i\)).

Use \({\bf 1}=(1,1)^t, Q=(1,-1)^t\), and unit columns by color. Let \(t_i=\tau_i\) times the corresponding unit column for holes, otherwise zero. For a lower gas object write \(p_i\) for its signed length, \(z_i=\sum q_\nu\theta_\nu\) for its virtual angular moment, \(\lambda_i\) for its color unit column; all three are zero on holes. Use addition of data for P. The following table collects the resulting five species in either color of sign \(q\). Here \(a_i\) means respectively \(\ell,d,h,h,d=h\), the jump charge is \(j_i={\bf1}^t(\lambda_i-t_i)\), and the external-drive coefficient is \(\xi_i=|p_i|-j_i\).
\[
\begin{array}{c|l|rrrr}
 \text{species}&\text{constituents}&a_i&p_i/q&j_i&\xi_i\\\hline
 \mathrm L&\text{one lower variable}&1/2&1&1&0\\
 \mathrm D&\text{two lower variables}&1/8&2&1&1\\
 \mathrm O&\text{outer hole}&3/8&0&-1&1\\
 \mathrm I&\text{inner hole}&-1/4&0&1&-1\\
 \mathrm P&\text{bound D and I}&0&2&2&0
\end{array}
\]
For fixed species, counts and positions, we will isolate all \(M\)-dependence in a Gaussian count factor. We first define its center and the residual interactions, state the resulting factorization, and then verify the contour and scalar bookkeeping.

Continue \(J_{\rm D}^\pm\) by retaining the initial branches (coefficient dependence \(e^{\pm d v}\)). Define columns
\[
\begin{gathered}
 F_i^+=J_i^++t_i e^{h_i v}=T_i/v+f_i+r_i v+\cdots,\\
 F_i^-=J_i^- -t_i e^{-h_i v}=-F_i^+(-v),\qquad
 T_i=Qp_i/4,\quad f_i=Qz_i/4-\lambda_i+t_i.
\end{gathered}
\]
(terms only on applicable constituents). The dressed oriented Fourier kernel \(G(v)\) is
\[
 G_{ij}=W^{g}_{ij}-(F_i^-)^t A^{-1}F_j^+
       +t_i^t t_j\,\operatorname{sgn}(h_i-h_j)e^{(h_j-h_i)v}+C_{ij}.
\]
Here \(W^g\) is just the preceding bare \(W\) between lower gas parts, with actual lower-lower shift branches. The crossing correction for each I,D pair with relative offsets \(h,d\) satisfying \(d>h\), writing \(R=1\) on equal colors and \(R=-2\) otherwise, is
\[
 C_{\rm ID}=R e^{(d-h)v},\qquad C_{\rm DI}=-R e^{(h-d)v}.
\]
Use half at equality, zero for \(d<h\). Expressions on constituent pairs are added. Use the winding \(k_{ij}\) from the direct interactions: gas winding already defined, plus \(-t_i^t t_j\,\operatorname{sgn}(h_i-h_j)\), plus \(-R\) from I to D in the crossing correction (antisymmetric, half at equality).

Write \(g\) for the column of lower virtual counts by color, \(u=M{\bf1}-g\). Thus \(M\ge\max g_c\) and the winding constraints for holes read
\[
 \sum(t_i-\lambda_i)=bQ .
\]
Let \(P_+,P_-\) be projections onto \({\bf1},Q\), \(a=9/8\), and put
\[
 \nu_i=-2T_i-P_+ r_i/a,\qquad
 M_0{\bf1}=n{\bf1}+\sum_i\nu_i+g .
\]
Indeed \(Q^t\nu_i=-p_i\), so the last equality is consistent. \(M_0=n+O(\#\mathrm{objects})\) is real and independent of positions and periods. The table's charges \(j=(1,1,-1,1,2)\) have sum zero over the configuration, by the winding constraint. At small \(v\),
\[
 G(v)=\frac{j j^t}{2av}+K_0+O(v),\qquad
 (K_0)_{ij}=\frac{p_i z_j-z_i p_j}{4}-k_{ij}
                    +\nu_i^t f_j-f_i^t\nu_j,\quad G(-v)=-G(v)^t.
\]
This follows by multiplying the columns using \(A=P_-/(2v)+A_1 v+O(v^3),\ P_+A_1=aP_+\); the bare pole cancels in the minus channel.

\paragraph{Fixed-sector reduction.}
For a fixed nonnegative \(M\) at large real \(L\), the termwise contour operations on the corresponding small-cycle sector of \(\mathcal N_n(e^{-L/2},s)\) give configurations of the five species above satisfying \(\sum_i j_i=0\), \(\sum_i(t_i-\lambda_i)=bQ\), and \(M\ge\max_c g_c\). Each term retains its exact one-object measures and factorials, the \(\kappa\)-free part of \(B_n\), and a constant phase of modulus one depending only on the discrete configuration. The factor \(\kappa^n\) from \(B_n\) combines with the length fugacity \(e^{-sM}\) to give the displayed count twist \(e^{-s(M-n)}\). The one-object prefactors are bounded by \(C|dx_i|\), uniformly for large \(\Re L\) and bounded \(\Im L\). These prefactors and the implicit phases are independent of \(M\) at fixed defect data. Apart from them, the term is the product of the following four factors; all remaining \(M\)-dependence occurs in the count factor.

\emph{Pressure and determinant.} The overall scalar factor is
\[
 \frac{2\pi}{L}\ \prod_{m>0}(\det A(m\varepsilon))^{-1}\,
 \exp\!\left[n^2\varepsilon\Big(\tfrac12 H_*(0)+\sum_{m>0} H_*(m\varepsilon)\Big)\right],
 \qquad H_*(v)=\frac{2\cosh3v-2}{v(2\cosh v-1)\sinh4v}.
\]
\emph{Count Gaussian and twist.} The count factor is
\[
 \exp[-2\varepsilon(b_1-b_2)^2-a\varepsilon(M-M_0)^2
              -s(M-n)+i(M-M_0)\mathcal S],\qquad
 \mathcal S=\sum_i j_i\phi_i.
\]
\emph{External drive.} Its logarithm is
\[
 -n\sum_i\xi_i\left[i\phi_i+a_i\varepsilon+
     \sum_{m>0}\frac{\rho(m\varepsilon)}{m}
             \big(e^{a_i m\varepsilon}y_i^m-e^{-a_i m\varepsilon}y_i^{-m}\big)\right],
 \qquad \rho(v)=\frac1{2\cosh v-1}.
\]
\emph{Interaction.} Its logarithm, in addition to a linear phase, is
\[
 -\frac\varepsilon2\sum_{ij} [v]G_{ij}
 -\sum_{m>0}\frac{\sum_{ij}G_{ij}(m\varepsilon)y_j^m/y_i^m-\sum_i(1-e^{-m\varepsilon})}{m}.
\]
Its additional linear phase is precisely \(-i\sum_{ij}\phi_j(K_0)_{ij}\).

The displayed count twist and linear phase retain their position dependence; only the implicit constant phases are suppressed. The verification below is termwise at fixed counts. Section~\ref{c01-c1-a-kernel-positivity-decomposition} bounds the interaction matrix frequency by frequency. Section~\ref{c01-c1-estimates-and-summation-of-the-defect-gas} then proves absolute summability, justifies the assembled fixed-sector identity, continues it to complex periods, and permits the sum over the original count.

\paragraph{Contour crossings and one-object factors.}
The growing term in the I-to-D \(- (F_i^-)^t A^{-1} F_j^+\) is \(-R e^{(d-h)v}\). Transferring the corresponding singular factor to the opposite Fourier direction adds \(C\) as above. Between two P objects at equal shifts both directions are transferred or averaged together (the powers are integer powers of the combined two-sided factor, up to a sign); on the internal I,D pair the singular factor is omitted by the residue.

Indeed at high positive \(v\), the only growing or nondecaying singular part of \(J_{\rm D}^+\) when crossing is the column with entries \(R e^{d v}\), whereas \(J_{\rm D}^-\) decays at least as fast as \(O(e^{-(1+d)v})\), and \(J_{\rm L}^\pm=O(e^{-v/2})\). This follows by summing the given divisors. Since \(A^{-1}=1+O(e^{-v})\), and D always stays at offset strictly inside O, the troublesome product among these factors is precisely \(-R e^{(d-h)v}\) from I to D, giving \((1-e^{(d-h)\varepsilon}y_D/y_I)^{-R}\); other than the direct hole pair and gas factors all the remaining mode sums are analytic across the moves. Re-expansion gives the correction, mean and phase above (log mean changes by \(-\varepsilon R(d-h)\)). The pole is simple and taking its residue removes this factor and the normalized I circle measure. Bound pairs have no further poles in these moves: equal-color DD and II each have a double zero at coincidence, compensating the possible simple pole with an unbound object or the two poles between bound pairs. Opposite colors had a zero rather than a pole; remaining bare factors between lower contours have no poles for these offset differences. Thus only disjoint matchings are taken, with factorial denominators now those for the combined species and signs depending on the defect data. This reasoning can move contours successively with separated radii. The natural expansions of lower-lower and hole-hole factors after each move can always be taken directly. Constants and phases for the continued upper-lower modes track the retained branches: under an offset increment \(\delta\), a coefficient continued by \(e^{\delta v}\) changes its mean \(p\alpha^2\theta^2/L-\varepsilon[v]W\) by exactly \(\varepsilon k\delta\). In paired-paired terms at equal offsets the two transferred factors have opposite orientations and can take the symmetric boundary values together up to signs.

More explicitly the leading terms of \(J_{\rm D}^+\) come from index \(-2\) at the top (same color) or \(-3\) at the bottom (opposite color), retained with wrapped value \(8+d\); the rest is \(O(e^{-(1-d)v})\). One can cross by moving the I contours while all D have \(d=-1/8\), and then move free D's and bound pairs (\(h=d\)) to final positions. Gas top relative offsets between two D gas parts then have size \(<1\) (the possible same-color constituent poles at zero top difference already cancel); differences from L to the two constituents have size \(<1\). Two P's on one contour combine their I-to-D pure factors to \(((1-y'/y)(1-y/y'))^{-R}\), agreeing with the averaged expansions. The gas windings group to integers (at equal D shifts the \(+1,-1\) constituent differences occur together), as do the crossing powers (the equal-offset halves occur in opposite directions together).

Self subtractions in the final series are \(G_{ii}(+\infty)(1-e^{-v})\) (always \(G_{ii}(+\infty)=1\)); self and measure prefactors besides powers already isolated are bounded by \(C\,|dx_i|\) per object, uniformly for large \(\Re L\) and bounded \(\Im L\). Indeed the lower self corrections were given earlier. An unbound hole has a missing 1 on the diagonal from the circle-pair log. Using \(1-e^{-v}\) instead leaves an activity \(\exp(\sum e^{-m\varepsilon}/m)\) with its measure \(dx_i/L\). For a bound pair, removing the pure pole adds 1 on its effective diagonal, so the two old subtractions (\(1-e^{-v}\) for D and 1 for I) now give exactly \(1-e^{-v}\) without extra activity. The residue has only the D radial measure. At infinity the joint kernel has same-color limit the identity in species, opposite-color limit 2 at P,P only, by the divisor rules.

\paragraph{Zero modes and count dependence.}
The direct constant of \(G\) at zero is \((p_i z_j-z_i p_j)/4-k_{ij}\) by symmetry of divisor shifts. Dressing contributes \(-2T_i^tf_j+2 f_i^tT_j\) and \((-r_i^tP_+f_j+f_i^tP_+r_j)/a\) there. For means during the contour moves, retaining a branch with offset increase \(\delta\) multiplies its coefficient by \(e^{\delta v}\); the winding-power increment of the constant log is then as above, using \(W^{[0]}=p\theta/4-k\) on a virtual pair. In total the \(z\)-square term precisely accounts for changes in all Gaussian offsets and the theta winding offsets by neutrality. Natural re-wrappings just use the actual direct factors (hole Cauchy, bare gas, or transferred pole).

Indeed the retained \([v]\) changes by \((p\theta/4-k)\delta+p\delta^2/8\) per virtual pair; the sum of \(p\) times squared shift differences changes by minus the change of \((\sum z_i)^2\). The winding monomials from the old formulas change by the same required scalar when an old radial argument is moved (\(y_{\rm old}=e^{\varepsilon\delta}y_{\rm new}\) for a single increment). On the internal bound pair \([v]C\) costs zero; in the nonzero modes deleting the pure ID pole subtracts \(1/m\) per frequency after identification, giving the self subtraction specified above.

For external insertions the dressed columns (rows similarly) couple with \(h_i^+-H{\bf1}^tA^{-1}F_i^+=\xi_i \rho e^{a_i v}\), reversed \(-\xi_i\rho e^{-a_i v}\); the exterior-exterior kernel is \(2H\rho=-vH_*\). Indeed \(-A^{-1}H{\bf1}=\rho {\bf1}\). On gas divisors, cyclic neighbor sum minus the divisor itself for \(h\) equals minus the sum of the two pair divisors. In multiplying \(h_i^+\) by \(2\cosh v-1\) the unwrapped neighbor exponents therefore leave only terms passing the seam. At the initial shift branches with \(d<0\) none do so for L or D bottom; at D top the \(-1\) index passes down through zero, leaving \(+e^{dv}\). This continues identically. For zero modes before dressing the bracket in the mean, now including the factors \(Y^{-\tau u_c}\), is
\[
 2|f_{\rm tot}|^2+4(b_1-b_2)^2
 +[v]\big\{u^tAu+2u^t(nH{\bf1}+\textstyle\sum F_i^+)+B(v)\big\}.
\]
Here \(f_{\rm tot}=(\sum z_i/4+b)Q\), the norm square is real, and \(B(v)\) denotes the sum of the direct kernels \(W^g,K^{\rm hole}=t_i^t t_j\operatorname{sgn}(h_i-h_j)e^{(h_j-h_i)v},C\) and direct external \(n(h_i^++h_i^-)\). Hole means and crossings indeed match these coefficients of \(v\). Replacing \(B(v)\) by
\[
 B(v)-(nH{\bf1}+\sum F_i^-)^t A^{-1}(nH{\bf1}+\sum F_i^+)
\]
incorporates all the other terms except \(4(b_1-b_2)^2+2a(M-M_0)^2\). This is completion of the square coefficientwise, using \(P_-u=-2\sum T_i, H=-av+O(v^3),\ \sum f_i=f_{\rm tot}\). Finally the winding phases on \(y_j\) have powers \(\sum_i k_{ij}+u^t(\lambda_j-t_j)-n|p_j|-bp_j\), with count signs independent of \(M\) as seen above. Splitting \(u\) by \(M_0\), the identity
\(\sum_i((K_0)_{ij}+k_{ij})=b p_j-\sum_i\nu_i^t(\lambda_j-t_j)\)
gives exactly the stated phases.

Explicitly in the square completion the minus-channel pieces with \(P_-u=-2T_{\rm tot}\) leave \(-2|f_{\rm tot}|^2\) at order \(v\), canceling the explicit norm square. The plus-channel part is \(a|P_+u+P_+ r_{\rm tot}/a-n{\bf1}|^2\) since \(P_+f_{\rm tot}=0\).

To reiterate the count-sign point, the upper cycle measure was \(-a_L\,dx/(2\pi i)=C_L e^{-s/2}dx/(2\pi)\) with exactly the positive-sign self conversion on real periods. The signs from \((-y_i)^{u_{c(i)}}\) over gas parts and \((-1)^{\sum u_c(e_c+bq_c)}\) from the winding determinant combine to a sign independent of \(M\) (\(e_c\) here counts those gas parts). Transferring pure factors and taking their normalized circle residues uses only signs fixed by defect data. These constant phase signs can simply be kept unchanged in continuation.

\subsection{A kernel positivity decomposition}\label{c01-c1-a-kernel-positivity-decomposition}

Our next task is to bound the residual particle sum uniformly in the spectral period. A positive Fourier form will control high local populations; a separate pole term will control the cumulative integer charge. The following decomposition supplies both controls.

In this calculation only the five species (per color) index the matrices; the symmetric part of \(G(v)\) for \(v>0\) is meant. One can subtract the following pieces:
\par\noindent\textbullet\quad \(j j^t\operatorname{sech}^2(Kv)/(2av)\), with \(K\) a sufficiently large constant, all colors;
\par\noindent\textbullet\quad on the opposite-color blocks only, the entrywise repulsion
\[
 R(v)=\tanh(v/2)\big(R_1/\cosh v+R_2/\cosh(v/2)\big)+2\tanh v\ e_5 e_5^t ,
\]
where
\[
 8 R_1=\begin{pmatrix}6&15&0&3&16\\15&22&0&9&26\\0&0&2&0&0\\3&9&0&0&5\\16&26&0&5&5\end{pmatrix},\qquad
 8 R_2=2(e_2e_5^t+e_5e_2^t)+5(e_4e_5^t+e_5e_4^t);
\]
\par\noindent\textbullet\quad on the diagonal, \(\tanh(c v)\), for sufficiently small \(c>0\).

The result divided by \(v\) is an even, smooth, exponentially decreasing positive Fourier form; it even dominates \(c_1\operatorname{sech}(c_2 v)\) times the identity, for some positive constants. The repulsions (second and third pieces divided by \(v\)) have entrywise nonnegative Fourier transforms. Indeed \(\tanh(c v)/v=\int_0^c\operatorname{sech}^2(u v)\,du\), and \(\operatorname{sech}\) has positive Fourier transform (by shifting through one pole), so squares and products here give positive convolutions.

Here are verification details of the matrix inequality. Work in either color channel \(h=\pm1\), i.e. the eigenblock same plus \(h\) times opposite, and first subtract only \(hR(v)\). Call the resulting symmetric matrix \(S_h(v)\). With \(x=e^{v/8}, w=x^8\), put \(D=w^4-w^{-4}\) and
\[
 A_n=w^{-4}(w^2-w+1)(w^3-h)^2.
\]
Thus the single upper eigenvalue is \(A_n/D\). After removal of the offsets (factor \(x^{8\beta},\ \beta=a_j-a_i\)), the oriented entry of \(G\) in this channel is
\[
 m_{ij}(w)/D - p_i(w^{-1})p_j(w)/(D A_n)+b_{ij}.
\]
Here \(m_{ij}\) is the divisor exponential sum for bare gas parts, \(p_i\) that from an upper variable plus \(\tau_i D\), and \(b_{ij}\) the hole-direct and crossing entries; explicit arrays are given below. At every \(w^8=1\), \(m_{ij}A_n=p_i(w^{-1})p_j(w)\), since wrapping and the branch choice do not change any powers. At \(w=1, h=1\) the difference vanishes to order at least two because \(p_i(1)=0\) and \(A_n\) has a double zero. Thus, putting \(z=w^{-1}\),
\[
 \begin{split}
 N&=44\ (h=1),\quad40\ (h=-1),\\
 \Delta_+&=(1-z)(1+z)^2(1+z^2)(1+z+z^2)^2(1-z+z^2),\\
 \Delta_-&=(1+z)^2(1+z^2)(1-z+z^2)^3 ,
 \end{split}
\]
the matrix \(4 x^N\Delta_h S_h\) is a Laurent polynomial of parity \(h\) under \(x\mapsto1/x\), with powers at most \(N\) (indeed \(S_h\to1\) at infinity). Denote its coefficients at \(x^k,\ k\ge0\) by \(P_k\), taking only half the constant coefficient. In the Taylor expansion in \(\log x\) the relevant matrices are
\[
 B_p=\sum_{k=0}^N (k/N)^p P_k
\]
for even \(p\) in the plus and odd \(p\) in the minus channel. \(B_0=128 j j^t\) for the plus channel (here \(j\) is the five-component vector). All the other relevant matrices are positive definite.

We include small algebra details for checking this last sign test. Take \(K_i=(2),(2,3),\varnothing,\varnothing,(2,3)\), respectively, and \(\tau_i=0,0,1,-1,-1\). Let \(d_l=(2,-1,1,-2h,2h,-2h,1,-1)_l\) for residues starting at 0. Put
\[
 L_b(m)=
 \begin{cases}
 z^{8-(m\bmod8)}&m\not\equiv0,\\
 z^8, 1, (1+z^8)/2&m\equiv0,\quad b>0, b<0, b=0\ \text{respectively}.
 \end{cases}
\]
Then
\[
 m'_{ij}=z^4 m_{ij}(w)=\sum_{r\in K_i,s\in K_j,l}d_l L_\beta(l+s-r),\quad
 u_j=z^4 p_j(w)=\tau_j(1-z^8)+\sum_{s\in K_j,l}d_l L_{\sigma_j}(l+s),
\]
where \(\sigma_j>0\) on L, negative on D or P. Finally
\[
 b_{ij}=-\tau_i\tau_j\operatorname{sgn}\beta+(1-2h)
 [\mathbf1_{\tau_i=-1,\,|K_j|=2}\chi(\beta)-
  \mathbf1_{|K_i|=2,\,\tau_j=-1}\chi(-\beta)]
\]
with \(\chi\) the positive step (half at zero).

The following table gives integer parts of the five successive leading minors; the last row in each block instead gives integer parts of \(100\sum_{k<N}\sum_j (k/N)^p |P_{k,ij}|\), by rows \(i\).
\[
\begin{array}{r|rrrrr@{\quad}r|rrrrr}
p(+)&1&2&3&4&5&p(-)&1&2&3&4&5\\\hline
2&22&149&542&2294&10192&1&1&11&9&9&3\\
4&10&32&31&54&18&3&1&9&9&11&9\\
6&7&15&28&97&174&5&1&12&14&17&21\\
8&5&13&38&135&211&7&2&15&22&31&59\\
10&4&12&44&150&177&9&2&17&29&43&122\\
12&4&13&47&158&144&11&3&18&33&52&185\\
14&4&13&49&165&125&13&3&17&36&63&238\\
16&4&14&50&173&122&15&3&17&38&77&289\\
18&4&14&52&180&133&17&3&16&40&96&347\\
20&4&14&53&188&156&19&3&16&43&115&412\\
22&4&15&54&196&188&21&3&16&45&135&482\\
24&4&15&56&203&227&23&3&16&48&153&550\\
26&4&15&56&209&271&25&3&15&50&169&614\\
28&4&15&57&215&318&27&3&15&52&183&672\\
30&4&15&58&220&366&29&261&374&287&271&231\\
32&4&15&59&225&414\\
34&4&15&59&228&461\\
36&4&15&60&232&506\\
38&4&15&60&235&549\\
40&4&15&60&237&589\\
42&178&295&150&177&382
\end{array}
\]
Since \(P_N=4I\), each last row proves strict diagonal dominance at its index. The factors \((k/N)^p\) for \(k<N\) decrease with \(p\), so diagonal dominance continues at every larger relevant index. The preceding rows prove positivity of the remaining \(B_p\) by their leading principal minors.

Here is one way of doing the additions/multiplications for the table. Form to order six inclusive
\[
 H_{ij}=\Delta_h(m'_{ij}+b_{ij})-Q_h(z) z^8 u_i(1/z)u_j(z),
 \quad Q_+=(1+z)^2(1+z^2)\sum_{d=0}^6 z^d,\quad Q_-=1+z^2 .
\]
For the two channel blocks separately, set
\[
 D_{l,ij}=2[z^{l/8+a_j-a_i}]H_{ij}
 -2h[z^{l/8}]\Delta_h\left(
 2\frac{1-z}{1+z}\left(\frac{z R_1}{1+z^2}+\frac{\sqrt z R_2}{1+z}\right)
 +2\frac{1-z^2}{1+z^2}e_5 e_5^t\right)_{ij}.
\]
A missing power means coefficient zero. Starting at \(l=0\) through \(N-1\), use \(D_l+D_l^t\) (so the first is \(4I\)). Indeed expansion of the displayed oriented entry gives \(H_{ij}\), since \((1-z^8)^{-1}\) can be dropped here. To obtain column \(b\) in a minor row, take
\[
 \left\lfloor
 \frac{\det[\sum_{l=0}^{N-1}(D_{l,ij}+D_{l,ji})(N-l)^p]_{i,j\le b}}{N^{pb}}
 \right\rfloor ;
\]
in the last row add \(|D_{l,bj}+D_{l,jb}|(N-l)^p\) for \(l>0\) instead and multiply by \(100/N^p\). These coefficient formulas use only polynomials of small degree.

For some further detail on the common denominators: at an eighth root \(w\) put \(U=\sum_l d_l w^{l-4}=\sum_l d_l w^{4-l}=A_n\). Then \(m_{ij}=U(\sum_{r\in K_i}w^{-r})(\sum_{s\in K_j}w^s)\), and \(p_i(w)=U\sum_{r\in K_i}w^r\) there. Once the \(D\) poles cancel this leaves the factor \((1-z+z^2)(1-hz^3)^2\), with at most a simple pole at \(z=1\) in the plus channel; the repulsion denominators use \((1+z)^2(1+z^2)\). Their combination thus divides \(\Delta_h\). The offsets only give integral powers of \(x\). Symmetrization is odd under \(v\mapsto-v\) before multiplication by \(4x^N\Delta_h\); this gives the stated parity. The plus constant Taylor value is \(2B_0\) and also follows immediately from \(\Delta_+\sim72v\) and the pole. For moments at positive index \(p\), only coefficients strictly above power zero in \(x\) matter, so \((N-1)/8+\max|a_i-a_j|\) controls the truncation in \(z\).

To pass from these coefficient tests to every \(v>0\), expand the Laurent polynomial in \(\log x=v/8>0\). Its parity and the half-constant convention give
\[
 \begin{split}
 4x^{44}\Delta_+S_+&=2\sum_{r\ge0}
       \frac{(44\log x)^{2r}}{(2r)!}B_{2r},\\
 4x^{40}\Delta_-S_-&=2\sum_{r\ge0}
       \frac{(40\log x)^{2r+1}}{(2r+1)!}B_{2r+1},
 \end{split}
\]
where each line uses its own channel's \(B_p\). All scalar weights are nonnegative, and \(\Delta_h>0\) for \(x>1\), proving that \(S_h\) is positive definite. Moreover \(B_1\) gives positivity of \(S_-/v\) at zero, while \(B_2\) gives the positive order-\(v\) coefficient of \(S_+\) on the orthogonal complement of its pole; the denominator corrections to \(B_0=128jj^t\) remain in the pole direction. Subtracting the indicated smoothed pole leaves a positive definite order-\(v\) coefficient for sufficiently large \(K\); away from zero this subtraction can be made arbitrarily small by further increasing \(K\). At infinity the result is still \(I+O(e^{-c_3 v})\), \(c_3>0\), so subtracting \(\tanh(c v)\) times identity for sufficiently small \(c\) is also allowed and leaves the asserted lower bound at all frequencies.

Indeed in the plus block the pole removed is \(j j^t\operatorname{sech}^2(Kv)/(a v)\), improving the regular linear coefficient in its pole direction arbitrarily as \(K\) increases. Once this works in a neighborhood of zero, increasing \(K\) only improves the pointwise form there. Similarly one can decrease \(c\) further on compact intervals after ensuring the inequalities near zero and infinity (at infinity first take \(2c<c_3\)).

\subsection{Estimates and summation of the defect gas}\label{c01-c1-estimates-and-summation-of-the-defect-gas}

The preceding matrix calculation is only a frequency-by-frequency statement. We must turn it into an absolutely summable particle expansion, continue it to complex periods, and then sum the original count \(M\). We do these operations in that order; the final complex-twist estimate will be obtained only after all three steps.

Take \(R=\Re L\) large, \(|\Im L|\) bounded, \(s\) in a compact subset of the strip \(|\Im s|<\pi\). Constants can depend on these strip bounds. Positions may run along a common contour \(\gamma(r)=r+i\chi(r)\), \(-R/2<r<R/2\), extended by \(\gamma(r+R)=\gamma(r)+L\), with \(\chi\) odd, bounded on the segment, piecewise smooth with Lipschitz constant \(\sigma\) sufficiently small. Split the parametrizing period into boxes of widths between 1 and 2 with populations \(n_d\); put \(m=\sum n_d\). We give details also justifying convergence of the expansions.

\paragraph{Population control.}
Isolate the smoothed Coulomb (pole) piece of the symmetric kernel as in the decomposition. Its coefficient of degree 1 at zero costs nothing in the mean term (sum of \(j\) vanishes). For the remaining symmetric kernel divided by \(v\), say \(f\), the interaction (zero mode included) uses the periodized line kernels
\[
 -\frac12\sum_{k\in\mathbb Z}\widehat f_{ij}((x_j-x_i+kL)/\alpha),\qquad
 \widehat f(t)=\int f(v)e^{itv}\,dv .
\]
On the diagonal species, the repulsion with numerator \(\tanh(cv)\) is taken here only between distinct objects; its self log with the \((1-e^{-v})\) subtraction is bounded per object, by summing (the difference vanishes at zero and decays at infinity). No opposite-color self term occurs.

We record the analytic bounds underlying this use of periodization. The regular functions involved are even, analytic near zero and in a thin strip about the real axis, and exponentially decreasing at both real ends also on tilted rays with sufficiently small slopes. The same holds for the functions giving symmetric repulsions except that they need only decrease as \(O(1/|v|)\). Indeed they have the displayed hyperbolic and exponential formulas, with poles only away from zero on the imaginary axis. In a small wedge \(|\Im t|\le\sigma |\Re t|\) their transforms are exponentially decreasing at infinity together with first derivatives; at zero the smooth rapidly decreasing pieces give bounded values and derivatives, the others need bounds \(O(1+|\log|t||)\) and \(O(1+1/|t|)\). To see this for \(\Re t>0\) deform the line into \(v=u+i(\delta+2\sigma_0|u|)\), small fixed \(\delta,\sigma_0>0\), \(\sigma\le\sigma_0\); similarly reflect for the other end. The integrals for real nonzero \(t\) allow the deformation also for decay of order 1 by oscillatory limits. Poisson summation in the form
\(\varepsilon\sum_{d\in\mathbb Z} f(d\varepsilon)e^{i d\varepsilon x/\alpha}=\sum_k \widehat f((x+kL)/\alpha)\)
then gives the kernels above, first for real periods and positions (for a logarithmic singularity by Fourier series off the diagonal, or smoothing by a narrow Gaussian first), then by continuation. Here the degree-1 mean of the original regular numerator supplies the half-weighted value of \(f(0)\).

On the real comparison line of period \(R\), the smooth positive form costs at least \(c_4\sum n_d^2\). Here is the passage from frequency positivity to this population bound. Subtract its diagonal minorant \(c_1\operatorname{sech}(c_2v)I\); the ordered point-measure quadratic form of the remainder, including its diagonal and zero mode, is nonnegative by the periodized Fourier series. The transform of the subtracted scalar kernel is nonnegative at every real distance and bounded below for two points in a box of width at most 2. If \(n_{d,t}\) counts type \(t\) in box \(d\), the retained energy is therefore at least
\[
 c\sum_d\sum_{t=1}^{10}n_{d,t}^2
 \ge \frac{c}{10}\sum_d n_d^2.
\]
This uses matrix positivity for the smooth remainder, not entrywise positivity of its transform. The separate repulsions between distinct objects are nonnegative in space, and the diagonal self correction costs only \(O(m)\). Changing to complex \(\gamma\) costs only \(C\sigma\sum n_d^2\): every lifted real difference changes by at most \(\sigma\) times its size, so the derivative bounds above give losses \(C\sigma\) times exponentially decreasing pair kernels. Their sums are bounded by the box square sum (Cauchy--Schwarz). Coincident positions can be disregarded. Thus these parts of the real log are bounded by \(-c_5\sum n_d^2+C m\).

For the antisymmetric part \(K=(G-G^t)/2\), even, smooth and exponentially decreasing, include the additional linear phase involving \(K_0\) in the interaction. This gives on each ordered pair a primitive, vanishing at difference zero, of
\[
 -\frac{i}{2\alpha}\sum_k \widehat K((x+kL)/\alpha).
\]
Indeed the linear phase inserts half the zero mode of the derivative of the sine sum. It is imaginary on the real comparison. On the complex curve its extra real part again costs only \(C\sigma\sum n_d^2\), applying the same endpoint comparisons to the primitives of each image term (subtract at \(kL\)), whose derivatives have the exponential bound. Thus small enough slopes allow absorption.

Here the lower endpoint comparison at \(k=0\) has exactly zero change, and those at nonzero \(k\) cost \(O(\sigma e^{-c'R})\) summed per pair, harmless also for the box bound. For the transform arguments with \(1/|v|\) tails, the nonnegative real transforms used above are locally integrable functions (use the displayed positive mixtures of scaled \(\operatorname{sech}^2\) transforms, with convolutions as applicable), with the pointwise bounds off zero; Poisson's formula there also follows by periodic Fourier coefficients in \(L^1\).

For example on a real period these series off coincidences converge in the ordinary paired-frequency sense too, since at positive frequencies one only has a constant over frequency plus a rapidly decreasing remainder at infinity. On straight complex periods, with \(x/L\) fixed real in the Poisson formula, both sides continue (for the image transforms keep the positive or negative real-end branch according to the lifted real difference). For the antisymmetric primitive, real parts can be compared image by image using a line primitive real on the real axis, evaluated at a lifted difference and at \(kL\), before multiplying by the imaginary prefactor. The joining paths can be chosen on the appropriate wedge side, since for differences inside one period the endpoint real parts do not straddle zero strictly.

\paragraph{Integer-charge variance and contour deformation.}
The population bound does not yet localize the objects along the period. For this we retain the pole term, whose charges are the integer jumps \(j_i\). Let \(w(r)\) on the real parametrizing period have initial value zero at \(-R/2\), with jumps \(j_i\) at the positions \(r_i\), periodically extended. Smooth this step function on the curve by convolution in contour length coordinate (complex differential \(dz\)) with
\[
 H(z)=\frac1{2\pi\alpha}\sum_k \widehat{\operatorname{sech}(K\,\cdot)}((z+kL)/\alpha)
\]
(here \(K\) in \(\operatorname{sech}\) is the pole smoothing constant). Denote the result by \(w_*\), and its average by \(\bar w_*\); this average is that of the steps using complex differential, \(-\mathcal S/(2\pi)\). The nonzero log of the smoothed pole is exactly
\[
 -c_0\int_\gamma(w_*-\bar w_*)^2 dz,\qquad c_0=\pi^2/(2\pi\alpha a).
\]
On straight contours these facts follow from Fourier series, differentiating the steps to the jump masses. On bent contours the identities continue: convolutions and integrals of steps use arcs with moving endpoints, the line kernels analytic with exponential bounds in neighborhoods of all the contour differences as contours and endpoints move together. Thus the resulting integrals continue holomorphically in endpoint data along the deformation (locally deform arcs by Cauchy's theorem). Equivalently the nonzero pole pair potential has second derivative given by a periodized smooth line kernel (numerator \(\operatorname{sech}^2\)) minus its constant mode.

In fact the unsmoothed step has nonzero Fourier coefficient \(\sum_i j_i e^{-i b\phi_i}/(2\pi i b)\) at integer \(b\ne0\), by integration by parts, and the smoothing multiplier is \(\operatorname{sech}(K b\varepsilon)\) (unit at zero).

We deform the product torus of spatial contours jointly. Three checks make this passage precise.

\emph{Coincidences.} Every lifted pair difference remains in the preceding wedge. Near a contact modulo $L$, separate the high-frequency constant of the \emph{combined} pair coefficient from its exponentially decreasing remainder. The former exponentiates to the nonnegative integer power $G_{ij}(+\infty)$ of $(1-y_j/y_i)(1-y_i/y_j)$; the remainder converges normally in a joint neighborhood. Thus the full exponential, rather than each separated logarithm, extends across contact loci. The external kernel used below has no singularity along this homotopy.

\emph{Period seams.} The original winding factors show periodicity: for each integer value of the original count $M$, the powers from winding and crossings are integers after constituents are combined. Continue the same straight-torus germ while following both representatives at a period seam. Their values remain equal. Small joint Cauchy deformations therefore create no vertical-end term. The common graph has a global Lipschitz bound on its imaginary part, with endpoint difference $\Im L$.

\emph{Moving endpoints.} For the step integrals, fix the order of distinct jump positions by real part and split every convolution arc at those positions. Expanding the convolutions gives iterated analytic-kernel integrals; all integration-point differences remain in the same analytic neighborhoods. Short tails added to the arcs show local holomorphy in their endpoints, including the smoothing variables. The arcs and their endpoints can consequently move together in small product neighborhoods. This continues the straight-contour identities needed at distinct jumps. These three checks justify the joint deformation without separating non-single-valued logarithmic terms.

Denote the smoothed step on the real comparison period by \(w_R\), and use \(b(r)\) here for a local density bound of the form \(\sum_{i,k} e^{-c_6|r-r_i+kR|}\), with \(c_6>0\) sufficiently small. The elementary smoothing estimates are
\[
 |w_R-w|\le C b(r),\qquad |w_*(\gamma(r))-w_R(r)|\le C\sigma b(r).
\]
Indeed both kernels integrate to one, so one convolves differences from the value at \(r\), bounded along a lift by the jumps passed. The kernels including differentials differ by at most a periodized exponential times \(C\sigma\). Thus \(\int b\le C m,\ \int b^2\le C\sum n_d^2\). Also, writing \(w_R^0\) for the real average, one has
\[
 |\bar w_* - w_R^0|\le C\sigma R^{-1/2}(\|w_R-w_R^0\|_2+\|b\|_2)
\]
by averaging the unsmoothed steps.

\paragraph{Sharp contour and absolute assembly.}
The external expression in brackets in the drive is a primitive (zero at 0) with derivative \((i/\alpha)\sum_k\widehat\rho((x+kL)/\alpha)\), evaluated at \(x=x_i-i\alpha a_i\). Here
\(\widehat\rho(t)\) is a positive constant times \(\sinh(2\pi t/3)/\sinh(\pi t)\), by shifting the integration line by \(2\pi i\) and taking the two poles. It decays as \(e^{-\pi |\Re t|/3}\), uniformly at large real ends, and for \(|\Im t|\le3/8\) has real part bounded below by a positive constant times the same decay (multiply by the conjugate denominator). Only D,O,I have nonzero drive. In general its real log on the contours is bounded above by \(C_n m\). For the sharp bound take
\[
 R=(3/2)\log n.
\]
This choice makes the weakest external drive, at the period edges, have order one: \(n e^{-2R/3}=1\). It also matches the pressure calculation in Section~\ref{c01-c1-bulk-calibration-and-puncture-exponent}: after removing its analytic correction \(c_*n^2e^{-L}\), the error is \(O(n^2e^{-4R/3})=O(1)\). The removed term itself has order \(\sqrt n\), so it cannot be included in the bounded error.

Bend so that \(\chi=0\) except within fixed width of the edge \(\{-R/2,R/2\}\), using sufficiently large constant width to achieve the small slope. At \(x_i=r_i\) the main line term integrated to \(r_i\) is imaginary, and on going vertically to the shifted position it gives signed real part at least \(c_7 e^{-4|r_i|/3}\), since \(\xi_i a_i>0\) for active objects. Other images contribute \(O(e^{-4(R-|r_i|)/3})\); the bend changes the result also by at most \(O(1/n)\) per primitive, by the fixed width and the decay. Hence the real log drive bound here is
\[
 C m-c_8\sum_{i:\xi_i\ne0}\exp(4\,\mathrm{dist}(r_i,\mathrm{edge})/3).
\]

Before resumming \(M\), these estimates, omitting common pressure and determinant factors, give (besides the \(-2\Re\varepsilon(b_1-b_2)^2\) log and factorial measures) bounds of the form
\[
 \exp[C_n m+C R-c_9\sum n_d^2-c_{10}(M-M_0)^2/R],
\]
where \(C_n\) can be taken constant in the sharp contour setting. Indeed the real part of the complex pole variance integral is bounded below (before its negative multiplier) by a constant times \(\|w_R-w_R^0\|_2^2-O(\sigma^2\|b\|_2^2)\), by the smoothing bounds. Also \(|\Im\mathcal S|\) is controlled by the complex average bound. Its cost multiplied by \(|M-M_0|\) is absorbed by a fraction of the count Gaussian and variance and box costs. The real tilt is bounded by completing a square, since \(M_0=n+O(m)\).

These bounds prove absolute summability after the stated transformations, even over all values of the original count \(M\) (the sum \(b_1+b_2\) is already determined by hole and lower counts). They upgrade the termwise reduction to an assembled fixed-sector identity. More explicitly, on the real periods one can first perform the circle expansion and shifts as above separately per fixed \(M,b_1,b_2\) and initial lower data, moving contours term by term at fixed hole counts; disjoint bindings just convert factorials by the matching choices. The bounds apply absolutely including these terms. On straight contours in a complex strip of large real periods all the expressions continue analytically with locally uniform convergence by the non-sharp bounds. Thus the identities for each original count \(M\) continue there, including vanishing after summing a sector \(M>2n\). We can bend terms as stated; their sum over \(M\) with only the lower constraint \(M\ge\max_c g_c\) is exactly \(\mathcal N_n(e^{-L/2},s)\). Extending this count sum to all integers is a subsequent approximation, with the same count-independent prefactors and negligible error in the sharp setting (using the unresummed bounds). If \(m\) is at most a small fraction of \(n\), the omitted Gaussian tail costs order \(n^2/R\); otherwise keep part of the box cost, at least a constant times \(m^2/R\). Unlocalized factorial measures even with \(\exp(C m)\) sum to \(\exp(O(R))\), so either error is negligible.

Indeed binding \(p\) indistinguishable pairs converts the factorials for the two counts to those for the residual counts and \(p!\); given the final counts they and the winding differences impose no new \(M\) constraint. Both error bounds here (relative to the common unsummed bulk factors) are negligible also against any \(\exp(-C R)\) on compact tilt ranges. Straight-contour holomorphy can first be used at fixed angular parametrizations, factoring the pure contact zeros and summing exponentially decreasing Fourier remainders locally uniformly for each fixed term, before continuing the torus as above.

\paragraph{Count summation and localization.}
Poisson summation in \(M\) now gives the prefactor \((\pi/(a\varepsilon))^{1/2}e^{-s(M_0-n)}\) and a sum over integers \(l\) with phases of modulus one (\(M_0\) real). Its exponential combined with the pole log is
\[
 \exp\left[-c_0\int_\gamma(w_*-l-is/(2\pi))^2 dz\right].
\]
Extract \(\exp(s^2/(4a\varepsilon))\). What is left costs in real log at least
\[
 c_{11}\|w_R-l\|_2^2-C m-C\sigma^2\|b\|_2^2
\]
(as suppression). Indeed the real linear coefficient relative to the square for real arguments, including the differential \(1+i\chi'\), has absolute value bounded strictly below 1, since \(|\Im s|<\pi\). For any such coefficient \(t\), \(x^2+t x\ge c_{12}x^2-C\,\mathrm{dist}(x,\mathbb Z)\) on the real line. The smoothed real step is within \(Cb\) of integer values, and replacing it by the complex smoothing costs only the controlled small errors above and a fraction of the square. Box costs again absorb quadratic errors.

Indeed (choosing the sign of \(l\)) the count transform's square exponent is \((s-i\mathcal S-2\pi i l)^2/(4a\varepsilon)=-c_0L(\bar w_*-l-is/(2\pi))^2\), completing the variance term. After the \(s^2\) extraction the real linear coefficient on replacing \(w_*\) by \(w_R\) is \(t=(\Im s+\chi'\Re s)/\pi\); slopes can be taken small depending on the prescribed compact tilt set. The additional replacement errors in the square and linear terms integrate against \(1+i\chi'\), and are bounded by a small fraction of \(\|w_R-l\|_2^2\) plus \(O(\sigma^2\|b\|_2^2+m)\).

Finally use \(g_0(r)=2\log(2+\mathrm{dist}(r,\mathrm{edge}))\). Integration by parts against the jumps gives
\[
 \sum j_i g_0(r_i)=-\int(w-l)g_0'\,dr
 \le\tfrac12 c_{11}\|w_R-l\|_2^2+C+C m .
\]
Only O has negative \(j\), and the external cost for O absorbs the wrong sign and leaves localization for it as well. Thus one keeps factors \(\exp(-g_0(r_i))\) on all objects, of bounded integral each. The unused variance penalty sums boundedly over \(l\). All defect sums with their bounded activities and factorials are now \(O(1)\). The Gaussian sum over \(b_1-b_2\) costs \(O(\sqrt R)\), offset together with the Poisson prefactor by the original \(2\pi/L\). This completes the bound, leaving only the bulk factors and the extracted harmonic twist.

\subsection{Bulk calibration and puncture exponent}\label{c01-c1-bulk-calibration-and-puncture-exponent}

The defect sum is bounded after extracting its explicit pressure, determinant and harmonic twist. We now compute these factors and compare them with Proposition~\ref{prop:pv:denominator}; this completes Theorem~\ref{thm:pv:puncture}.

For the overall determinant, Euler summation gives
\[
 -\sum_{m>0}\log\det A(m\varepsilon)=L/4+O(1).
\]
Indeed the eigenvalues are \((2\cosh v-1)(\cosh3v\mp1)/\sinh4v\); the negative log determinant extends smoothly at zero, decays exponentially, and its integral on the positive line is \(\pi^2/8\), by integrating
\(2[\log(1-e^{-8v})-\log(1-e^{-6v})-\log(1-e^{-v}+e^{-2v})]\).
Bounded Euler error follows also on the slightly rotated rays by the derivative bound there.

Put \(I=\int_0^\infty H_*(v)dv\). Poisson summation for the even extension shows
\[
 \varepsilon\left(\tfrac12 H_*(0)+\sum_{m>0}H_*(m\varepsilon)\right)
 =I+c_* e^{-L}+O(e^{-4\Re L/3})
\]
for an absolute constant \(c_*\). The transform at positive frequencies \(2\pi k/\varepsilon\) can use the same tilted line shift as above, here passing the first two simple poles at \(i\pi/4,i\pi/3\) (removable poles, if any, simply give zero), which gives this estimate including bounded imaginary parts of the period. For the bulk integral use \(y=e^{-v}\):
\[
 v H_*(v)=\frac{2y^2(1-y^3)^2}{(1-y+y^2)(1-y^8)}
          =\frac{P(y)}{1-y^{24}}
\]
with symmetric coefficients \(P_j=P_{24-j}\), sum zero; at \(j<12\) the nonzero ones are \(P_2=P_3=2,\ P_5=P_6=-6,\ P_8=P_9=8,\ P_{10}=2,P_{11}=-6\). Integrate groupwise against \(dv/v\) (absolute convergence by the double zero of \(P\)). Pairing the resulting log products symmetrically about 12 gives
\[
 I=-\sum_{j<12} P_j\log\sin(\pi j/24),\qquad
 \exp I=8D_0^3/81.
\]
Here one uses \(\prod_{l\ge0}(l+x)(l+1-x)/(l+1/2)^2=\sin(\pi x)\); e.g. log differentiation gives \(\pi\cot(\pi x)\) by its elementary partial fractions (residue integration on expanding half-integer squares), and the normalization is at \(1/2\). To simplify, the sines at 2 and 10 have product \(1/4\), at 5 and 11 product \((1+\sqrt2)/4\), at 3 and 9 product \(\sqrt2/4\).

Multiply the analytic numerator \(\mathcal N_n(\mathfrak q,s)\) by \(\exp(-c_* n^2\mathfrak q^2)\). On \(|\mathfrak q|=n^{-3/4}\) the bounds proved above apply around the entire circle, with the pressure error after multiplication bounded in log. The correcting factor equals one at zero, so Cauchy's formula gives
\[
 \mathcal N_n(0,s)=\frac{1}{2\pi i}
 \int_{|\mathfrak q|=n^{-3/4}}
 e^{-c_*n^2\mathfrak q^2}\mathcal N_n(\mathfrak q,s)
 \frac{d\mathfrak q}{\mathfrak q}.
\]
This bounds the original numerator \(S_{2n}(1)^2X_{n,n}(\kappa=e^s)\). The \(\kappa\)-free part of \(B_n\) times the bulk growth is exactly \(D_0^{-2n}(4D_0^2/9)^{2n^2}\), the squared exponential factor in the denominator bound. The remaining factors besides constants are bounded by \(\left|\exp[L/4+s^2/(4a\varepsilon)]\right|\). Since \(a=9/8\), \(\varepsilon=2\pi\alpha/L\), and \(\alpha=\pi/4\), their powers on this circle satisfy, uniformly on each compact tilt set,
\[
 \left|\exp[L/4+s^2/(4a\varepsilon)]\right|
 \le C n^{3/8+2\Re(s^2)/(3\pi^2)},\qquad
 \frac38-2\cdot\frac5{48}=\frac16.
\]
Thus division by the squared homogeneous Pfaffian gives, locally uniformly in \(|\Im s|<\pi\),
\[
 |E_n(s)|\le C,\quad
 E_n(s)=X_{n,n}(\kappa=e^s)\,\exp[-(1/6+2s^2/(3\pi^2))\log n].
\]
At \(s=i\pi/2\) this analytic function is exactly 1. Hence all these puncture polynomials as functions of \(\eta=2\cosh s\) are zero-free in a fixed disk about \(\eta=0\), with Taylor coefficients of each fixed order in \(\log X_{n,n}\) bounded by \(O_k(1+\log n)\) (Cauchy estimates for \(\log E_n\) near the normalization point). Also
\[
X_{n,n}(1)=n^{1/6+o(1)}
\]
where 1 denotes \(\kappa=1\), thus fugacity 2. For the lower bound take any subsequence and a further locally uniform limit of \(E_n\) in the strip. It is nonzero at the normalization point and hence nonzero also at points \(s=i t\) with \(t>0\) arbitrarily small. There \(2\cos t<2\), so coefficient positivity gives the lower bound with exponent \(1/6-2t^2/(3\pi^2)\). The upper bound uses \(s=0\) directly.

Here compact subsequential convergence follows e.g. from derivative bounds by Cauchy's estimates followed by a compact exhaustion and diagonal selection, with holomorphy of the limit. The same bounds at \(i\pi/2\) give a neighborhood where \(|E_n-1|<1/2\) uniformly, and \(\eta\) has a local analytic inverse parameter there since its derivative is nonzero.
